% ============================================================
% EN3150 - ASSIGNMENT 03
% RESOURCE-CONSTRAINED CNN FOR EDGE IMAGE CLASSIFICATION
% ============================================================

\documentclass[12pt,a4paper]{report}

% ============================================================
% PACKAGES
% ============================================================

\usepackage[a4paper,
            left=1in,
            right=1in,
            top=1in,
            bottom=1in]{geometry}

\usepackage{graphicx}
\usepackage{float}
\usepackage{booktabs}
\usepackage{array}
\usepackage{tabularx}
\usepackage{longtable}
\usepackage{multirow}
\usepackage{amsmath}
\usepackage{amssymb}
\usepackage{enumitem}
\usepackage{caption}
\usepackage{subcaption}
\usepackage{setspace}
\usepackage{titlesec}
\usepackage{tocloft}
\usepackage{fancyhdr}
\usepackage{url}
\usepackage{xcolor}
\usepackage{hyperref}
\usepackage{background}

% ============================================================
% GENERAL SETTINGS
% ============================================================

\onehalfspacing

\setlength{\parindent}{0pt}
\setlength{\parskip}{0.4em}

%----------------------------------------------------------------------------------------
% PAGE BORDER
%----------------------------------------------------------------------------------------

\backgroundsetup{
  scale=1,
  color=black,
  opacity=1,
  angle=0,
  contents={%
    \begin{tikzpicture}[remember picture, overlay]
      \draw[thick]
      ([xshift=1cm, yshift=-1cm]current page.north west)
      rectangle
      ([xshift=-1cm, yshift=1cm]current page.south east);
    \end{tikzpicture}
  }
}
% ============================================================
% HYPERREF
% ============================================================

\hypersetup{
    colorlinks=false,
    hidelinks,
    pdftitle={Resource-Constrained CNN for Edge Image Classification},
    pdfauthor={Dilhara D.S.},
    pdfsubject={EN3150 Assignment 03}
}

% ============================================================
% CHAPTER FORMAT
% ============================================================

\titleformat{\chapter}
    {\normalfont\bfseries\Large}
    {\thechapter}
    {1em}
    {}

\titlespacing*{\chapter}
    {0pt}
    {0pt}
    {18pt}

% ============================================================
% SECTION FORMAT
% ============================================================

\titleformat{\section}
    {\normalfont\bfseries\large}
    {\thesection}
    {0.8em}
    {}

\titleformat{\subsection}
    {\normalfont\bfseries\normalsize}
    {\thesubsection}
    {0.8em}
    {}

\titlespacing*{\section}
    {0pt}
    {10pt}
    {5pt}

\titlespacing*{\subsection}
    {0pt}
    {8pt}
    {4pt}

% ============================================================
% TABLE OF CONTENTS FORMAT
% ============================================================

\renewcommand{\contentsname}{Contents}

\renewcommand{\cftchapfont}{\bfseries}
\renewcommand{\cftchapleader}{\cftdotfill{\cftdotsep}}

\renewcommand{\cftsecleader}{\cftdotfill{\cftdotsep}}

\renewcommand{\cftsubsecleader}{\cftdotfill{\cftdotsep}}

\setlength{\cftbeforechapskip}{5pt}
\setlength{\cftbeforesecskip}{1pt}
\setlength{\cftbeforesubsecskip}{0pt}

% ============================================================
% FIGURE CAPTION
% ============================================================

\captionsetup{
    font=small,
    labelfont=bf,
    justification=centering
}

% ============================================================
% HEADER AND FOOTER
% ============================================================

\pagestyle{fancy}

\fancyhf{}

\fancyfoot[C]{\thepage}

\renewcommand{\headrulewidth}{0pt}
\renewcommand{\footrulewidth}{0pt}

% Chapter pages also use same footer
\fancypagestyle{plain}{
    \fancyhf{}
    \fancyfoot[C]{\thepage}
    \renewcommand{\headrulewidth}{0pt}
    \renewcommand{\footrulewidth}{0pt}
}

% ============================================================
% HELPER COMMAND FOR OPTIONAL FIGURES
% ============================================================

\newcommand{\insertfigure}[4]{
    \begin{figure}[H]
        \centering

        \IfFileExists{#1}{
            \includegraphics[
                width=#2\textwidth
            ]{#1}
        }{
            \fbox{
                \parbox[c][5cm][c]{0.80\textwidth}{
                    \centering
                    Figure file not found:\\
                    \texttt{#1}
                }
            }
        }

        \caption{#3}
        \label{#4}
    \end{figure}
}

% ============================================================
% DOCUMENT START
% ============================================================

\begin{document}

% ============================================================
% TITLE PAGE
% ============================================================

\begin{titlepage}

\centering

\vspace*{0.3cm}

% ------------------------------------------------------------
% UNIVERSITY
% ------------------------------------------------------------

{\Huge\bfseries
University of Moratuwa
\par}

\vspace{0.25cm}

{\Large\bfseries
Faculty of Engineering
\par}

\vspace{0.8cm}

% ------------------------------------------------------------
% UNIVERSITY LOGO
% ------------------------------------------------------------

\IfFileExists{figures/uom_logo.png}{
    \includegraphics[
        width=0.30\textwidth
    ]{figures/uom_logo.png}
}{
    \fbox{
        \parbox[c][4cm][c]{0.30\textwidth}{
            \centering
            University Logo
        }
    }
}

\vspace{1cm}

% ------------------------------------------------------------
% COURSE
% ------------------------------------------------------------

{\Large\bfseries
EN3150 -- Pattern Recognition
\par}

\vspace{0.7cm}

{\LARGE\bfseries
Assignment 03
\par}

\vspace{0.5cm}

{\Large\bfseries
Resource-Constrained CNN for
Edge Image Classification
\par}

\vspace{0.5cm}



% ------------------------------------------------------------
% GROUP INFORMATION
% ------------------------------------------------------------

{\Large\bfseries
Group Members
\par}

\vspace{0.4cm}

\begin{center}
\renewcommand{\arraystretch}{1.3}

\begin{tabular}{|p{8cm}|p{4cm}|}
\hline
\textbf{Name} & \textbf{Index Number} \\
\hline
Dilhara D.S. & 230147L \\
\hline
 DISSANAYAKE G.R.G.K.  & 230159B \\
\hline
MENDIS T.A.S.D.  & 230409T \\
\hline
 MARASINGHA M.D.S.  & 230398F \\
\hline

\end{tabular}

\end{center}

\vfill

{\large
2026
\par}

\end{titlepage}

% ============================================================
% TABLE OF CONTENTS
% ============================================================

\pagenumbering{roman}

\tableofcontents

\clearpage

% ============================================================
% MAIN CONTENT
% ============================================================

\pagenumbering{arabic}

% ============================================================
% CHAPTER 1
% ============================================================

\chapter{Dataset Overview: MNIST 2-Digit Dataset}

% ------------------------------------------------------------
\section{Overview}
% ------------------------------------------------------------

The dataset used in this project is the MNIST 2-Digit Dataset.
The dataset contains handwritten images representing two-digit
numbers.

Unlike the original MNIST dataset, which contains ten individual
digit classes, the two-digit dataset represents numbers using two
digits. Therefore, the classification task contains 100 possible
classes, ranging from 00 to 99.

The objective of this assignment is to investigate convolutional
neural network architectures for two-digit handwritten number
classification while considering memory and computational
constraints.

The project evaluates both custom CNN architectures and
fine-tuned pre-trained models.

% ------------------------------------------------------------
\section{Characteristics of the Dataset}
% ------------------------------------------------------------

The dataset provides image files together with CSV metadata.

The metadata contains two digit labels:

\begin{itemize}
    \item \texttt{label\_1}: First digit.
    \item \texttt{label\_2}: Second digit.
\end{itemize}

The final class label is calculated as:

\[
\text{Class Label}
=
10\times\text{label\_1}
+
\text{label\_2}
\]

Therefore:

\[
00,01,02,\ldots,98,99
\]

represent the 100 possible classes.
\newpage
% ------------------------------------------------------------
\section{Image Pre-processing}
% ------------------------------------------------------------

The assignment specifies a maximum image resolution of
$64\times64$ pixels.

Therefore, all images were resized to:

\[
64\times64
\]

The preprocessing pipeline consisted of:

\begin{enumerate}
    \item Loading image metadata.
    \item Constructing the two-digit class label.
    \item Locating the corresponding image files.
    \item Resizing images to $64\times64$ pixels.
    \item Converting images into tensors.
    \item Normalizing the image data.
\end{enumerate}

% ------------------------------------------------------------
\section{Train-Validation-Test Split}
% ------------------------------------------------------------

The dataset was divided into three subsets:

\begin{itemize}
    \item 70\% training data
    \item 15\% validation data
    \item 15\% test data
\end{itemize}

A stratified split was used so that the class distribution was
maintained as consistently as possible across the training,
validation, and testing subsets.
\newpage
% ------------------------------------------------------------
\section{Sample Images}
% ------------------------------------------------------------

Figure~\ref{fig:figures/dataset_samples} shows examples of handwritten
two-digit images from the dataset.

\insertfigure
    {figures/dataset_samples.png}
    {0.90}
    {Sample images from the MNIST 2-Digit Dataset.}
    {fig:dataset_samples}

% ------------------------------------------------------------
\section{Dataset Summary}
% ------------------------------------------------------------

\begin{table}[H]
\centering
\caption{Summary of the MNIST 2-Digit Dataset}
\label{tab:dataset_summary}

\begin{tabular}{ll}
\toprule
\textbf{Property} & \textbf{Value} \\
\midrule
Dataset & MNIST 2-Digit Dataset \\
Task & Two-digit number classification \\
Number of Classes & 100 \\
Class Range & 00--99 \\
Maximum Image Resolution & $64\times64$ \\
Training Split & 70\% \\
Validation Split & 15\% \\
Testing Split & 15\% \\
\bottomrule
\end{tabular}

\end{table}


% ============================================================
% CHAPTER 2
% ============================================================

\chapter{Custom CNN for Handwritten Digit Classification}

% ------------------------------------------------------------
\section{Dataset Loading}
% ------------------------------------------------------------

The dataset metadata was loaded from the provided CSV files.

The CSV files contain the image filename and the two individual
digit labels.

The image filename was matched recursively with the corresponding
image file.

The two labels were combined into one 100-class target label.

% ------------------------------------------------------------
\section{Dataset Pre-Processing}
% ------------------------------------------------------------

Each image was converted into a grayscale image and resized to
$64\times64$ pixels.

The preprocessing pipeline was kept consistent between the
custom CNN models.

This ensures that differences in performance are primarily caused
by architectural differences rather than different input
processing procedures.

% ------------------------------------------------------------
\section{Train-Validation-Test Split}
% ------------------------------------------------------------

The same 70/15/15 split was used for all custom models.

This is important because the models must be evaluated using
the same test distribution to provide a fair comparison.

% ------------------------------------------------------------
\section{CNN Model Architecture}
% ------------------------------------------------------------

Two custom CNN architectures were developed.

\begin{itemize}
    \item Model A: Standard CNN.
    \item Model B: Lightweight CNN using depthwise-separable
    convolutions.
\end{itemize}

% ------------------------------------------------------------
\subsection{Model A Architecture}
% ------------------------------------------------------------

Model A is a standard convolutional neural network consisting
of three convolutional blocks followed by fully connected layers.
\newpage
The architecture is:

\begin{enumerate}
    \item Conv2D: $1\rightarrow16$
    \item ReLU
    \item MaxPooling
    \item Conv2D: $16\rightarrow32$
    \item ReLU
    \item MaxPooling
    \item Conv2D: $32\rightarrow64$
    \item ReLU
    \item MaxPooling
    \item Fully Connected: $64\times8\times8\rightarrow128$
    \item ReLU
    \item Fully Connected: $128\rightarrow100$
\end{enumerate}

The final layer produces 100 class scores.

\begin{figure}[H]
    \centering
    \includegraphics[width=1\linewidth]{figures/Model A.png}
    \caption{Model A Architecture}
    \label{fig:placeholder}
\end{figure}

% ------------------------------------------------------------
\subsection{Model A Architecture Summary}
% ------------------------------------------------------------

\begin{table}[H]
\centering
\caption{Model A Architecture}
\label{tab:model_a_arch}

\begin{tabular}{lll}
\toprule
\textbf{Layer} & \textbf{Output Channels} & \textbf{Operation} \\
\midrule
Input & 1 & Grayscale image \\
Conv1 & 16 & $3\times3$ convolution \\
Pool1 & 16 & Max pooling \\
Conv2 & 32 & $3\times3$ convolution \\
Pool2 & 32 & Max pooling \\
Conv3 & 64 & $3\times3$ convolution \\
Pool3 & 64 & Max pooling \\
FC1 & 128 & Fully connected \\
Output & 100 & Classification \\
\bottomrule
\end{tabular}

\end{table}

% ------------------------------------------------------------
\section{Justification for Activation Functions}
% ------------------------------------------------------------

The ReLU activation function was used throughout the custom CNN.

The ReLU function is:

\[
f(x)=\max(0,x)
\]

ReLU is computationally simple and introduces non-linearity into
the network.

It also avoids some of the optimization difficulties associated
with saturating activation functions.

% ------------------------------------------------------------
\section{Why Adam Optimizer is Used?}
% ------------------------------------------------------------

Adam was selected as the primary optimizer for the main CNN
experiments.

Adam combines adaptive learning rates with momentum-based
optimization.

This makes it suitable for training CNNs efficiently without
requiring extensive manual learning-rate adjustment.

% ------------------------------------------------------------
\section{Learning Rate Selection}
% ------------------------------------------------------------

The learning rate controls the size of the parameter updates
during training.

A learning rate that is too large may result in unstable
optimization, while a very small learning rate may result in
slow convergence.

The selected learning rate was used consistently for the main
model comparison.

% ------------------------------------------------------------
\section{CNN Model Training}
% ------------------------------------------------------------

The custom CNN models were trained for 20 epochs.

Cross-entropy loss was used for the 100-class classification
problem.

The training process monitored:

\begin{itemize}
    \item Training loss
    \item Training accuracy
    \item Validation loss
    \item Validation accuracy
\end{itemize}

% ------------------------------------------------------------
\section{Model A Training Results}
% ------------------------------------------------------------

Model A achieved strong performance during training.

The final test results are shown in Table~\ref{tab:model_a_results}.

\begin{table}[H]
\centering
\caption{Model A Test Performance}
\label{tab:model_a_results}

\begin{tabular}{lr}
\toprule
\textbf{Metric} & \textbf{Result} \\
\midrule
Trainable Parameters & 560,612 \\
Model Size & 2.1431 MB \\
Average Time/Epoch & 52.27 s \\
Test Accuracy & 94.06\% \\
Macro Precision & 93.91\% \\
Macro Recall & 94.94\% \\
\bottomrule
\end{tabular}

\end{table}

% ------------------------------------------------------------
\section{Model A Learning Curves}
% ------------------------------------------------------------

\insertfigure
    {figures/model_a_curves.png}
    {0.85}
    {Losses of Training and validation curves for Model A.}
    {fig:figures/model_a_curves}
\insertfigure
    {figures/model_a_curves_.png}
    {0.85}
    {Accuracy Training and validation curves for Model A.}
    {fig:figures/model_a_curves_}

% ------------------------------------------------------------
\newpage
\section{Optimizer Comparison}
% ------------------------------------------------------------

Three optimizers were investigated:

\begin{enumerate}
    \item Adam
    \item SGD
    \item SGD with Momentum
\end{enumerate}

A five-epoch comparison was conducted to study their early
optimization behaviour.

% ------------------------------------------------------------
\subsection{Comparison Metrics}
% ------------------------------------------------------------

The main metrics used for the optimizer comparison were:

\begin{itemize}
    \item Training accuracy
    \item Validation accuracy
    \item Training loss
    \item Validation loss
\end{itemize}

% ------------------------------------------------------------
\subsection{Optimizer Results}
% ------------------------------------------------------------

The optimizer comparison produced the following results.

\begin{table}[H]
\centering
\caption{Five-Epoch Optimizer Comparison}
\label{tab:optimizer_comparison}

\begin{tabular}{lcc}
\toprule
\textbf{Optimizer} &
\textbf{Best Validation Accuracy} &
\textbf{Final Validation Accuracy} \\
\midrule
SGD & 4.89\% & 4.90\% \\
SGD + Momentum &  14.86\% & 5.62\% \\
Adam & 10.15\% & 7.92\% \\
\bottomrule
\end{tabular}

\end{table}

The SGD with Momentum configuration reached a peak validation
accuracy of approximately 21.36\% during the short experiment.

However, its final validation accuracy dropped to approximately
7.97\%.

Adam reached higher training accuracy during the five-epoch
experiment, reaching approximately 8.41\% training accuracy by
epoch 5.

% ------------------------------------------------------------
\section{Impact of Momentum Parameter}
% ------------------------------------------------------------

Momentum influences the direction and magnitude of parameter
updates.

The general momentum update can be represented as:

\[
v_t=\beta v_{t-1}+(1-\beta)g_t
\]

where:

\begin{itemize}
    \item $v_t$ is the current velocity,
    \item $\beta$ is the momentum coefficient,
    \item $g_t$ is the current gradient.
\end{itemize}

Momentum can improve convergence, but inappropriate settings can
also cause unstable optimization.

% ------------------------------------------------------------
\subsection{Comparison of Momentum Behaviour}
% ------------------------------------------------------------

The results indicate that momentum can produce rapid improvements
during the early training stage.

However, the validation accuracy did not remain stable throughout
the short five-epoch experiment.

Therefore, a longer experiment would be required before making a
final conclusion about the best momentum value.

% ------------------------------------------------------------
\subsection{Key Observations}
% ------------------------------------------------------------

The main observations are:

\begin{itemize}
    \item Optimizer choice strongly affects convergence.
    \item Momentum can improve early learning.
    \item Validation performance can fluctuate significantly.
    \item Five epochs are insufficient for a definitive optimizer
    benchmark.
\end{itemize}

% ------------------------------------------------------------
\section{Model Evaluation on Test Dataset}
% ------------------------------------------------------------

After training, the final models were evaluated using the
independent test dataset.

The following metrics were calculated:

\begin{itemize}
    \item Accuracy
    \item Macro precision
    \item Macro recall
    \item Confusion matrix
\end{itemize}

% ------------------------------------------------------------
\subsection{Overall Performance Metrics}
% ------------------------------------------------------------

Model A achieved:

\[
Accuracy=94.06\%
\]
\[
Precision=93.91\%
\]
\[
Recall=94.94\%
\]

These results indicate strong classification performance across
the 100 classes.

% ------------------------------------------------------------
\subsection{Confusion Matrix Analysis}
% ------------------------------------------------------------

\insertfigure
    {figures/model_a_confusion_matrix.png}
    {0.80}
    {Confusion matrix of Model A on the test dataset.}
    {fig:figures/model_a_confusion}

% ------------------------------------------------------------
\subsection{Analysis and Insights}
% ------------------------------------------------------------

The high accuracy, precision, and recall indicate that Model A
successfully learned useful visual features from the handwritten
two-digit images.

The model provides a good balance between model complexity and
classification performance.


% ============================================================
% CHAPTER 3
% ============================================================

\chapter{Resource-Constrained CNN}

% ------------------------------------------------------------
\section{Motivation for a Lightweight Model}
% ------------------------------------------------------------

Edge devices often have limited memory, processing power, and
energy availability.

Therefore, a model with a large number of parameters may not be
appropriate for deployment on resource-constrained hardware.

The assignment requires a custom CNN with fewer than 100,000
trainable parameters.

Model B was therefore designed specifically to satisfy this
constraint.

% ------------------------------------------------------------
\section{Depthwise-Separable Convolution}
% ------------------------------------------------------------

Model B uses depthwise-separable convolutions.

A depthwise-separable convolution divides a conventional
convolution into:

\begin{enumerate}
    \item Depthwise convolution.
    \item Pointwise convolution.
\end{enumerate}

This reduces the number of parameters and computational operations
compared with standard convolution.

% ------------------------------------------------------------
\section{Model B Architecture}
% ------------------------------------------------------------

The architecture consists of:

\begin{enumerate}
    \item Depthwise-separable convolution:
    $1\rightarrow16$
    \item Max pooling
    \item Depthwise-separable convolution:
    $16\rightarrow32$
    \item Max pooling
    \item Depthwise-separable convolution:
    $32\rightarrow64$
    \item Max pooling
    \item Adaptive average pooling
    \item Fully connected layer: $64\rightarrow100$
\end{enumerate}

\begin{figure}[H]
    \centering
    \includegraphics[width=1\linewidth]{figures/Model B.png}
    \caption{Lightweight Model}
    \label{fig:placeholder}
\end{figure}
% ------------------------------------------------------------
\section{Parameter Constraint Verification}
% ------------------------------------------------------------

The assignment requires:
\[
\text{Trainable Parameters}<100,000
\]

Model B contains:
\[
9,741
\]
trainable parameters.

Therefore:
\[
9,741 < 100,000
\]

The model successfully satisfies the required resource constraint.

% ------------------------------------------------------------
\section{Model A Learning Curves}
% ------------------------------------------------------------

\insertfigure
    {figures/model_b_curves.png}
    {0.85}
    {Losses of Training and validation curves for Model B.}
    {fig:figures/model_b_curves}
\insertfigure
    {figures/model_b_curves_.png}
    {0.85}
    {Accuracy Training and validation curves for Model B.}
    {fig:figures/model_b_curves_}
    
% ------------------------------------------------------------
\section{Model B Evaluation}
% ------------------------------------------------------------

The final Model B results are shown below.

\begin{table}[H]
\centering
\caption{Model B Test Performance}
\label{tab:model_b_results}

\begin{tabular}{lr}
\toprule
\textbf{Metric} & \textbf{Result} \\
\midrule
Trainable Parameters & 9,741 \\
Model Size & 0.0466 MB \\
Average Time/Epoch & 50.69 s \\
Test Accuracy & 62.64\% \\
Macro Precision & 51.69\% \\
Macro Recall & 47.87\% \\
\bottomrule
\end{tabular}

\end{table}

% ------------------------------------------------------------
\section{Model B Training Behaviour}
% ------------------------------------------------------------

The training accuracy increased to approximately 77.49\% by
epoch 20.

However, the final validation accuracy decreased to approximately
23.84\% at epoch 20.

This indicates instability near the end of training and suggests
that further regularization or learning-rate tuning may improve
generalization.

% ------------------------------------------------------------
\section{Model B Confusion Matrix}
% ------------------------------------------------------------

\insertfigure
    {figures/model_b_confusion_matrix.png}
    {0.80}
    {Confusion matrix of Model B on the test dataset.}
    {fig:figures/model_b_confusion}

% ------------------------------------------------------------
\section{Model B Analysis}
% ------------------------------------------------------------

Model B provides a very small memory footprint.

It contains only 9,741 trainable parameters compared with
560,612 parameters in Model A.

Therefore, Model B has approximately 57.6 times fewer trainable
parameters than Model A.

However, this large reduction in model complexity resulted in a
substantial reduction in classification accuracy.

Model B achieved 64.97\% accuracy compared with 94.35\% for
Model A.

This demonstrates the trade-off between resource usage and
classification performance.


% ============================================================
% CHAPTER 4
% ============================================================

\chapter{Comparison with State-of-the-Art Networks (Transfer Learning)}

% ------------------------------------------------------------
\section{Selection and Justification of Pre-trained Models}
% ------------------------------------------------------------

Two established convolutional neural network architectures were
selected for transfer learning:

\begin{enumerate}
    \item EfficientNet-B0
    \item MobileNetV2
\end{enumerate}

These architectures were selected because they are widely used
for efficient image classification and provide a useful
comparison with the custom resource-constrained architecture.

% ------------------------------------------------------------
\subsection{Model 1: EfficientNet-B0}
% ------------------------------------------------------------

EfficientNet-B0 uses compound scaling to balance network depth,
width, and input resolution.

It provides strong classification performance while maintaining
a relatively efficient architecture compared with many larger
CNNs.

% ------------------------------------------------------------
\subsection{Model 2: MobileNetV2}
% ------------------------------------------------------------

MobileNetV2 was designed specifically for efficient computation
on mobile and embedded devices.

It uses inverted residual blocks and linear bottlenecks to reduce
computational cost.

% ------------------------------------------------------------
\section{Fine-Tuning Strategy and Implementation}
% ------------------------------------------------------------

Both pre-trained models were adapted for the 100-class
two-digit classification problem.

The final classification layer was replaced with a new layer
producing 100 output classes.

% ------------------------------------------------------------
\subsection{Data Preprocessing Pipeline}
% ------------------------------------------------------------

The custom dataset images are grayscale.

The pre-trained networks expect three-channel input.

Therefore, the grayscale images were converted into a
three-channel representation before being passed to the
pre-trained networks.

All models used the same underlying train, validation, and test
data split.

% ------------------------------------------------------------
\subsection{Model Adaptation Methodology}
% ------------------------------------------------------------

The classifier layer of each pre-trained network was replaced
with a 100-class output layer.

The models were then fine-tuned using the project dataset.

The objective was to determine whether transfer learning could
provide higher classification accuracy than the custom CNNs.

% ------------------------------------------------------------
\subsection{Transfer Learning Approach}
% ------------------------------------------------------------

The pre-trained ImageNet weights provide an initial set of
learned visual features.

Fine-tuning then adapts these representations to the
two-digit handwritten classification problem.

This approach provides a useful comparison between:

\begin{itemize}
    \item Learning features completely from scratch.
    \item Reusing features learned from a large external dataset.
\end{itemize}

% ------------------------------------------------------------
\section{Individual Model Performance Evaluation}
% ------------------------------------------------------------

The two pre-trained models were trained for 20 epochs and
evaluated on the same test set.

% ------------------------------------------------------------
\subsection{MobileNetV2 Fine-tuned Model Evaluation}
% ------------------------------------------------------------

The final MobileNetV2 results are:

\begin{table}[H]
\centering
\caption{MobileNetV2 Performance}
\label{tab:mobilenet_results}

\begin{tabular}{lr}
\toprule
\textbf{Metric} & \textbf{Result} \\
\midrule
Trainable Parameters & 2,351,972 \\
Model Size & 9.2100 MB \\
Average Time/Epoch & 99.50 s \\
Test Accuracy & 97.14\% \\
Macro Precision & 97.04\% \\
Macro Recall & 97.56\% \\
\bottomrule
\end{tabular}

\end{table}

MobileNetV2 achieved higher accuracy than both custom CNN
architectures.

\insertfigure
    {figures/MobileNetV2.png}
    {0.85}
    {Losses of Training and validation curves for MobileNetV2.}
    {fig:figures/MobileNetV2}
\insertfigure
    {figures/MobileNetV2_.png}
    {0.85}
    {Accuracy of Training and validation curves for MobileNetV2.}
    {fig:figures/MobileNetV2_}
% ------------------------------------------------------------
\subsection{EfficientNet-B0 Fine-tuned Model Evaluation}
% ------------------------------------------------------------

The final EfficientNet-B0 results are:

\begin{table}[H]
\centering
\caption{EfficientNet-B0 Performance}
\label{tab:efficientnet_results}

\begin{tabular}{lr}
\toprule
\textbf{Metric} & \textbf{Result} \\
\midrule
Trainable Parameters & 4,135,648 \\
Model Size & 16.0674 MB \\
Average Time/Epoch & 89.48 s \\
Test Accuracy & 97.82\% \\
Macro Precision & 97.60\% \\
Macro Recall & 98.10\% \\
\bottomrule
\end{tabular}

\end{table}

EfficientNet-B0 achieved the highest test accuracy among the
four evaluated models.

\insertfigure
    {figures/EfficientNet-B0.png}
    {0.85}
    {Losses of Training and validation curves forEfficientNet-B0.}
    {fig:figures/EfficientNet-B0}
\insertfigure
    {figures/EfficientNet-B0_.png}
    {0.85}
    {Accuracy of Training and validation curves for EfficientNet-B0.}
    {fig:figures/EfficientNet-B0_}
% ------------------------------------------------------------
\section{Performance Comparison: Custom CNN vs. Fine-tuned
State-of-the-Art Models}
% ------------------------------------------------------------

The complete comparison is shown in Table~\ref{tab:complete_comparison}.

\begin{table}[H]
\centering
\caption{Comprehensive Performance Comparison}
\label{tab:complete_comparison}

\resizebox{\textwidth}{!}{
\begin{tabular}{lrrrrrr}
\toprule
\textbf{Model}
&
\textbf{Parameters}
&
\textbf{Size (MB)}
&
\textbf{Time/Epoch}
&
\textbf{Accuracy}
&
\textbf{Precision}
&
\textbf{Recall}
\\
\midrule

Model A
&
560,612
&
2.1431
&
77.48 s
&
94.06\%
&
93.91\%
&
94.94\%
\\

Model B
&
9,741
&
0.0466
&
69.69 s
&
62.64\%
&
51.69\%
&
47.87\%
\\

MobileNetV2
&
2,351,972
&
9.2100
&
101.95 s
&
97.14\%
&
97.04\%
&
97.56\%
\\

EfficientNet-B0
&
4,135,648
&
16.0674
&
118.78 s
&
97.82\%
&
97.60\%
&
98.10\%
\\

\bottomrule
\end{tabular}
}

\end{table}

% ------------------------------------------------------------
\section{Point-by-Point Comparative Analysis}
% ------------------------------------------------------------

\subsection{Accuracy}

EfficientNet-B0 achieved the highest test accuracy:

\[
97.82\%
\]

MobileNetV2 achieved:

\[
97.14\%
\]

Model A achieved:

\[
94.06\%
\]

Model B achieved:

\[
62.82\%
\]

Therefore, according to accuracy:

\[
\text{EfficientNet-B0}
>
\text{MobileNetV2}
>
\text{Model A}
>
\text{Model B}
\]
\begin{figure}[H]
    \centering
    \includegraphics[width=1\linewidth]{figures/accuracy_vs_model.png}
    \caption{Accuracy vs Model}
    \label{fig:placeholder}
\end{figure}
% ------------------------------------------------------------
\subsection{Memory Footprint}
% ------------------------------------------------------------

Model B has the smallest model size:

\[
0.0466\text{ MB}
\]

This is significantly smaller than all other evaluated models.

EfficientNet-B0 has the largest model size:

\[
16.0674\text{ MB}
\]

% ------------------------------------------------------------
\subsection{Computational Cost}
% ------------------------------------------------------------

Model A has the shortest average training time per epoch among
the four models:

\[
77.48\text{ seconds}
\]

Model B takes:

\[
69.69\text{ seconds}
\]

Despite having substantially fewer parameters, Model B is not
faster than Model A.

This demonstrates that parameter count alone does not determine
actual execution time.

Implementation efficiency, memory access, hardware utilization,
and convolution implementation can all affect runtime.

% ------------------------------------------------------------
\section{Key Findings and Interpretations}
% ------------------------------------------------------------

The main findings are:

\begin{enumerate}
    \item EfficientNet-B0 achieved the highest accuracy.
    \item Model B achieved the smallest memory footprint.
    \item Model A provides a strong middle ground.
    \item MobileNetV2 provides high accuracy with a smaller
    footprint than EfficientNet-B0.
    \item Parameter reduction can result in substantial accuracy
    degradation.
    \item A smaller parameter count does not necessarily result
    in faster training.
\end{enumerate}

% ------------------------------------------------------------
\section{Strategic Implications}
% ------------------------------------------------------------

The choice of model should depend on the requirements of the
deployment environment.

For an application where maximum classification accuracy is the
main objective, EfficientNet-B0 is the strongest model evaluated.

For extremely memory-constrained edge devices, Model B may be
preferred despite its lower accuracy.

Model A provides a useful compromise between model size,
training speed, and classification performance.

% ------------------------------------------------------------
\section{Conclusion}
% ------------------------------------------------------------

The experiments demonstrate that model architecture selection
requires balancing accuracy, memory footprint, and computational
cost.

The highest-performing model is not necessarily the most
appropriate model for every deployment environment.

EfficientNet-B0 achieved the highest accuracy, while Model B
achieved the lowest memory footprint.

Therefore, model selection should be based on the actual
constraints and requirements of the target application.

% ------------------------------------------------------------
\section{Trade-offs, Advantages, and Limitations:
Custom Model vs. Pre-trained Model}
% ------------------------------------------------------------

% ------------------------------------------------------------
\subsection{Custom CNN Model Analysis}
% ------------------------------------------------------------

\textbf{Advantages}

\begin{itemize}
    \item Full control over the architecture.
    \item Can be specifically designed for the target problem.
    \item Model A has a relatively small footprint compared with
    the pre-trained networks.
    \item No dependency on ImageNet feature representations.
    \item Architecture can be optimized specifically for edge
    deployment.
\end{itemize}

\textbf{Limitations}

\begin{itemize}
    \item Requires manual architecture design.
    \item Hyperparameter tuning may be necessary.
    \item Extremely small models can lose significant accuracy.
    \item Generalization to other datasets may be limited.
\end{itemize}

% ------------------------------------------------------------
\subsection{Pre-trained Model Analysis}
% ------------------------------------------------------------

\textbf{Advantages}

\begin{itemize}
    \item Strong feature representations.
    \item High classification accuracy.
    \item Established architectures.
    \item Useful for rapid experimentation.
    \item Efficient transfer learning when sufficient resources
    are available.
\end{itemize}

\textbf{Limitations}

\begin{itemize}
    \item Larger parameter count.
    \item Larger memory footprint.
    \item Higher computational cost.
    \item Additional input adaptation may be required.
    \item Pre-trained features may not be perfectly aligned with
    handwritten digit data.
\end{itemize}

% ------------------------------------------------------------
\subsection{Decision Framework: When to Choose Which Approach}
% ------------------------------------------------------------

\begin{table}[H]
\centering
\caption{Model Selection Guidelines}
\label{tab:model_selection}

\begin{tabularx}{\textwidth}{p{4cm}p{3.5cm}X}
\toprule
\textbf{Scenario}
&
\textbf{Recommended Model}
&
\textbf{Reason}
\\
\midrule

Maximum classification accuracy
&
EfficientNet-B0
&
Highest test accuracy among evaluated models.
\\

High accuracy with moderate resources
&
Model A
&
Good balance between accuracy, size, and training time.
\\

Extremely limited memory
&
Model B
&
Only 9,741 trainable parameters and 0.0466 MB model size.
\\

High accuracy with transfer learning
&
MobileNetV2
&
Strong accuracy with smaller footprint than EfficientNet-B0.
\\

Strict sub-100k parameter requirement
&
Model B
&
Successfully satisfies the assignment constraint.
\\

\bottomrule
\end{tabularx}

\end{table}

% ------------------------------------------------------------
\subsection{Strategic Trade-offs Summary}
% ------------------------------------------------------------

\textbf{Accuracy vs. Memory}

Higher accuracy was generally associated with larger models in
this experiment.

Model B demonstrates the extreme memory-efficient end of the
design space, while EfficientNet-B0 provides the highest accuracy.

\textbf{Efficiency vs. Accuracy}

Model A provides a practical compromise.

It achieved 94.06\% accuracy with only 2.1431 MB of model storage.

\textbf{Parameter Count vs. Runtime}

Model B contains dramatically fewer parameters than Model A,
but its average training time is slightly higher.

Therefore, parameter count should not be used as the only measure
of computational efficiency.

% ------------------------------------------------------------
\subsection{Conclusion: Balanced Strategy}
% ------------------------------------------------------------

The experiments demonstrate that there is no universally optimal
architecture.

EfficientNet-B0 is the best choice when classification accuracy
is the highest priority and sufficient computational resources
are available.

Model B is the most suitable option when memory is the dominant
constraint.

Model A provides a practical middle ground between the two
extremes.

Therefore, the final model should be selected based on the target
deployment requirements rather than accuracy alone.


% ============================================================
% CHAPTER 5
% ============================================================

\chapter{Overall Conclusion}

This project investigated resource-constrained convolutional
neural networks for classifying two-digit handwritten numbers.

Four architectures were evaluated:

\begin{enumerate}
    \item Standard custom CNN (Model A).
    \item Lightweight depthwise-separable CNN (Model B).
    \item MobileNetV2.
    \item EfficientNet-B0.
\end{enumerate}

Model A achieved 94.35\% test accuracy using 560,612 trainable
parameters.

Model B successfully satisfied the assignment requirement of
fewer than 100,000 trainable parameters, containing only 9,741
parameters. However, its test accuracy was 64.97\%.

MobileNetV2 achieved 96.79\% test accuracy.

EfficientNet-B0 achieved the highest accuracy at 97.47\%.

The results demonstrate an important trade-off between
classification performance and resource requirements.

For maximum accuracy, EfficientNet-B0 is the strongest choice.

For extremely constrained memory environments, Model B provides
the smallest model footprint.

Model A provides a balanced alternative with strong accuracy
and relatively low resource requirements.

Overall, the experiment demonstrates that designing neural
networks for edge devices requires considering accuracy,
memory footprint, and computational cost simultaneously.


% ============================================================
% REFERENCES
% ============================================================

\clearpage

\begin{thebibliography}{99}

\bibitem{lecun1998}
Y. LeCun, L. Bottou, Y. Bengio, and P. Haffner,
``Gradient-based learning applied to document recognition,''
\textit{Proceedings of the IEEE},
vol. 86, no. 11, pp. 2278--2324, 1998.

\bibitem{howard2017}
A. G. Howard et al.,
``MobileNets: Efficient Convolutional Neural Networks for Mobile
Vision Applications,''
arXiv:1704.04861, 2017.

\bibitem{sandler2018}
M. Sandler, A. Howard, M. Zhu, A. Zhmoginov, and L.-C. Chen,
``MobileNetV2: Inverted Residuals and Linear Bottlenecks,''
Proceedings of the IEEE Conference on Computer Vision and Pattern
Recognition, 2018.

\bibitem{tan2019}
M. Tan and Q. V. Le,
``EfficientNet: Rethinking Model Scaling for Convolutional Neural
Networks,''
Proceedings of the International Conference on Machine Learning,
2019.

\bibitem{kingma2014}
D. P. Kingma and J. Ba,
``Adam: A Method for Stochastic Optimization,''
arXiv:1412.6980, 2014.

\bibitem{he2016}
K. He, X. Zhang, S. Ren, and J. Sun,
``Deep Residual Learning for Image Recognition,''
Proceedings of the IEEE Conference on Computer Vision and Pattern
Recognition, 2016.

\bibitem{pytorch}
PyTorch Documentation,
``PyTorch: An Open Source Machine Learning Framework,''
Available at:
\url{https://pytorch.org/}

\bibitem{torchvision}
Torchvision Documentation,
``Torchvision Models and Datasets,''
Available at:
\url{https://pytorch.org/vision/}

\bibitem{transfer}
PyTorch Tutorials,
``Transfer Learning for Computer Vision Tutorial,''
Available at:
\url{https://pytorch.org/tutorials/beginner/transfer_learning_tutorial.html}

\end{thebibliography}


% ============================================================
% APPENDIX
% ============================================================

\appendix

\chapter{Appendix}

% ------------------------------------------------------------
\section{GitHub URL}
% ------------------------------------------------------------

The source code, notebooks, experiment results, and supporting
files are available in the following GitHub repository:

\begin{center}

\url{https://github.com/sachi1645/EN3150-Two-Digit-MNIST-Classification}

\end{center}



% ------------------------------------------------------------
\section{Final Model Comparison}
% ------------------------------------------------------------

\begin{table}[H]
\centering
\caption{Final Summary of All Evaluated Models}
\label{tab:final_summary}

\begin{tabular}{lrrrr}
\toprule
\textbf{Model}
&
\textbf{Parameters}
&
\textbf{Size}
&
\textbf{Time/Epoch}
&
\textbf{Accuracy}
\\
\midrule

Model A
&
560,612
&
2.1431 MB
&
77.48 s
&
94.06\%
\\

Model B
&
9,741
&
0.0466 MB
&
69.69 s
&
62.64\%
\\

MobileNetV2
&
2,351,972
&
9.2100 MB
&
101.95 s
&
97.14\%
\\

EfficientNet-B0
&
4,135,648
&
16.0674 MB
&
118.78 s
&
97.82\%
\\

\bottomrule
\end{tabular}

\end{table}



% ============================================================
% END DOCUMENT
% ============================================================

\end{document}